% TOP_PROBLEM: {"id": 3373, "title": "Are there only finite Fermat Primes?", "category_id": 1, "category": {"id": 1, "name": "number_theory", "display_name": "Number Theory", "description": "Properties of integers, prime numbers, Diophantine equations.", "slug": "number-theory", "order_index": 1, "created_at": "2026-07-31T15:26:25.670Z"}, "status": "open"}
% FIRST_POSTED: 2026-09-27
% !TeX program = lualatex
% ATTEMPT_STATUS: unresolved
\documentclass[11pt]{article}
\usepackage[margin=25mm]{geometry}
\usepackage{fontspec}
\setmainfont{Latin Modern Roman}
\usepackage{amsmath,amssymb,mathtools}
\usepackage{unicode-math}
\setmathfont{Latin Modern Math}
\usepackage{xurl,hyperref}
\hypersetup{colorlinks=true,urlcolor=blue}
\setcounter{secnumdepth}{0}
\setlength{\parindent}{0pt}
\setlength{\parskip}{5pt}
\setlength{\emergencystretch}{3em}
\begin{document}
\section*{\texorpdfstring{Are there only finite Fermat Primes?}{Are there only finite Fermat Primes?}}\label{3373}
\subsection{Definitions and mathematical statement}
For each integer $n\ge0$, the Fermat number is
\[
 F_n=2^{2^n}+1.
\]
A Fermat prime is a member of this sequence with no positive divisors other than one and itself. Determine whether the index set
\[
 \{n\in{\mathbb{Z}}_{\ge0}:F_n\text{ is prime}\}
\]
is infinite or finite. This is a binary determination problem, not a requirement to prove infinitude in advance. Proving compositeness for any finite set of further terms does not establish eventual compositeness, and hence does not by itself establish finiteness of the prime terms.
\par\smallskip\textit{Formulation and concept references:} \href{https://webusers.imj-prg.fr/~michel.waldschmidt/articles/pdf/NTchallengesVI.pdf}{[S1]}.

\subsection{Short English statement}
Are there infinitely many prime Fermat numbers, or are only finitely many terms of this specific sequence prime?
\subsection{Research attempt}
\paragraph{Scope and literature check (27 September 2026).}
This GPT-6 Astra Ultra attempt treats the question as a choice between
infinitely many and finitely many prime terms. Neither alternative is
assumed in advance. The elementary facts about the first five terms and
Euler's factor of the sixth are recorded in [S1]. Tsumura's primary
paper [E1] gives other primality tests and recalls the classical setting.
A test for one input is not a theorem about infinitely many indices.
The arguments here rederive the relevant certificates and investigate
two possible uniform approaches: divisibility obstructions and the
iteration underlying P\'epin's criterion.

No global resolution was established by this source check. In particular,
the numerical frontier in an old survey is not treated as a current
exhaustive computation. Only small terms actually verified below are
used in the finite record. The partial results are the exact prime-divisor
restriction, a limitation on every fixed finite sieve, and a complete
primality certificate whose quantifiers remain finite.

\paragraph{Lemma 476.1: the product identity gives new divisors, not new prime terms.}
For every $n\geq1$,
\begin{equation}\label{a476:product}
                  \prod_{j=0}^{n-1}F_j=F_n-2.
\end{equation}
For $n=1$ both sides are three. Multiplying the identity by
$F_n=2^{2^n}+1$ gives
\[
 (2^{2^n}-1)(2^{2^n}+1)=2^{2^{n+1}}-1=F_{n+1}-2,
\]
which proves the induction. If $m<n$, a common divisor of $F_m$ and
$F_n$ divides both $F_n-2$ and $F_n$, hence divides two. Every Fermat
number is odd, so
\begin{equation}\label{a476:coprime}
                         \gcd(F_m,F_n)=1\quad(m\ne n).
\end{equation}
Consequently there are infinitely many different primes among the prime
divisors of all the $F_n$: choose one divisor from each term, and none
can repeat. This conclusion allows every sufficiently late $F_n$ to be
composite. Pairwise coprimality does not say that any chosen term has
only one prime factor, or has exponent one in its factorization.

The distinction can be made without relying on the conjecture. List the
ordinary primes $p_1,p_2,\ldots$ and put $a_j=p_{2j-1}p_{2j}$.
Then the $a_j$ are pairwise coprime and all composite. This is not a
model of the Fermat recurrence, but it disproves the inference from
pairwise coprimality alone to infinitely many prime terms.

\paragraph{Lemma 476.2: exact order restrictions on a prime divisor.}
If an odd prime $q$ divides $F_n$, then
\begin{equation}\label{a476:order}
                         \operatorname{ord}_q(2)=2^{n+1}.
\end{equation}
Indeed $2^{2^n}\equiv-1\pmod q$, so the order divides $2^{n+1}$
and does not divide $2^n$. Every divisor of $2^{n+1}$ is a power of
two; the order must therefore be the full power. Lagrange's theorem in
$(\mathbb Z/q\mathbb Z)^\times$ gives $2^{n+1}\mid q-1$.

For $n\geq2$ there is one further factor of two:
\begin{equation}\label{a476:factorform}
                       q\equiv1\pmod{2^{n+2}}.
\end{equation}
The first restriction already gives $q\equiv1\pmod8$. The quadratic
character formula $(2/q)=(-1)^{(q^2-1)/8}$ then makes two a quadratic
residue. Euler's criterion gives $2^{(q-1)/2}\equiv1\pmod q$, so
the order $2^{n+1}$ divides $(q-1)/2$. This proves
\eqref{a476:factorform}. These familiar quadratic-reciprocity facts
are standard arithmetic inputs; no probabilistic assumption is involved.
The exceptions matter: $F_0=3$ and $F_1=5$ do not obey the strengthened
modulus written with their respective indices.

Conversely, a prime has order $2^{n+1}$ for two if and only if it
divides $F_n$. If the order is that number, $2^{2^n}$ has order two
modulo $q$, and the only element of order two in the prime field is
$-1$. Thus the prime-divisor set is described exactly by an order
condition, not just by an arithmetic progression. Most primes in the
progression need not meet the order condition.

For example $641=5\cdot2^7+1$ is of the required form for $n=5$.
To certify its divisibility without a large decimal division, observe
\[
                   641=5^4+2^4=5\cdot2^7+1.
\]
Modulo $641$, raising $5\cdot2^7\equiv-1$ to the fourth power and
substituting $5^4\equiv-2^4$ gives $-2^{32}\equiv1$. Hence
$641\mid F_5$, and it is a proper factor since $1<641<F_5$.
In fact $F_5=641\cdot6700417$, an equality checked by integer
multiplication. The argument for compositeness does not require a
primality proof for either factor.

\paragraph{A precise obstruction to a fixed finite congruence cover.}
Suppose one proposes to prove eventual compositeness using finitely many
fixed odd primes $q_1,\ldots,q_r$, by showing that every sufficiently
large $F_n$ is divisible by one of them. This is impossible. By
\eqref{a476:coprime}, each $q_i$ divides at most one member of the
entire sequence, so the proposed set covers at most $r$ indices.

More quantitatively, for any fixed bound $B$, at most $\pi(B)$ Fermat
numbers have a prime divisor at most $B$. Thus all sufficiently late
terms have no prime divisor below $B$. This statement is compatible
both with large prime terms and with products of large primes. It does
not give a single uniform bound on how many prime factors occur.

In fact \eqref{a476:factorform} supplies the stronger explicit lower
bound $q\geq2^{n+2}+1$ for $n\geq2$. Since
$\log_2F_n$ is about $2^n$, such a lower bound is still far smaller
than $\sqrt{F_n}$. It leaves ample room for proper factors. A test
looking only for small divisors eventually finds none for structural
reasons, whether or not the input is prime.

This excludes a particular common proof strategy, not every finite
description of an infinite divisibility argument. A rule producing a
different prime or composite divisor depending on $n$ would evade the
obstruction and could be decisive, but no such rule for all large $n$
is provided here. Dirichlet's theorem about primes in a progression
also does not supply a divisor of a prescribed integer in that progression.

\paragraph{Proposition 476.3: P\'epin's criterion, with both directions proved.}
For every $n\geq1$,
\begin{equation}\label{a476:pepin}
 F_n\text{ is prime}
 \quad\Longleftrightarrow\quad
                    3^{(F_n-1)/2}\equiv-1\pmod{F_n}.
\end{equation}
\emph{Proof of necessity.} Put $p=F_n$ and assume it is prime.
Since $2^n$ is positive and even, $2^{2^n}\equiv4\pmod{12}$, so
$p\equiv5\pmod{12}$. The quadratic-character formula for three
gives $(3/p)=-1$ in this residue class. Euler's criterion is precisely
the displayed congruence.

\emph{Proof of sufficiency.} Suppose the congruence holds and let $q$
be any prime divisor of $F_n$. Reduction modulo $q$ gives the same
congruence. In particular $q\ne3$. Write $F_n-1=2^M$ with $M=2^n$.
The order of three modulo $q$ divides $2^M$, since the square of
$3^{2^{M-1}}$ is one, but it does not divide $2^{M-1}$, since
$-1\ne1$ modulo the odd prime $q$. Hence its order is $2^M=F_n-1$.
It divides $q-1$, giving $q\geq F_n$. Since $q\mid F_n$, necessarily
$q=F_n$, proving primality. This also excludes composite prime powers;
no squarefreeness was assumed.\hfill$\square$

The case $n=0$ is handled separately: $F_0=3$ is prime but the base
three congruence is not a valid test for it. A Fermat congruence with
base two would be useless for the present purpose: every $F_n$, prime
or composite, already satisfies $2^{2^n}\equiv-1$. The special base
and exponent in \eqref{a476:pepin} are doing real work.

\paragraph{The test as an exact finite computation.}
Set $r_0=3$ and repeatedly square modulo $F_n$:
\[
                       r_{j+1}\equiv r_j^2\pmod{F_n}.
\]
After exactly $2^n-1$ squarings, $r_j$ represents
$3^{2^{2^n-1}}=3^{(F_n-1)/2}$. Testing whether the final residue is
$F_n-1$ therefore decides the primality of that one term. The number
of squarings is exponential in the index $n$, but only linear in the
binary length of $F_n$, which is $2^n+1$. This separates the short
description of the input from the length of the integer actually processed.

The companion exact-integer calculation checks $1\leq n\leq10$.
It certifies primality for $n=1,2,3,4$ and rejects all the tested
$n=5,\ldots,10$. Together with the separate check of $F_0$, this
recovers the five small prime terms. Every rejected input is composite
by the proved equivalence, whether or not a factor is found. These
eleven finite cases do not certify a cutoff after which all terms are
composite.

An inductive shortcut would need to relate the final residues at two
successive indices. The underlying moduli do satisfy
\begin{equation}\label{a476:recurrence}
                     F_{n+1}=(F_n-1)^2+1,
\end{equation}
but they are coprime. A congruence modulo $F_n$ does not determine a
congruence modulo $F_{n+1}$. Indeed the Chinese remainder theorem
allows any residue modulo the first modulus to coexist with any residue
modulo the second for a suitable integer. This observation is not a
claim that the particular powers of three are independent; it identifies
why a proof must use more than reduction of a previous residue.

\paragraph{The heuristic track and its exact logical limit.}
A very crude random-integer heuristic assigns an integer near $F_n$
a prime probability of order $1/\log F_n\asymp2^{-n}$. The sum over
$n$ converges. If one actually had random events with summable
probabilities, the first Borel--Cantelli lemma would imply that only
finitely many occur almost surely. Fermat primality is a deterministic
sequence, and no probability space with a theorem transferring that
conclusion to this sequence has been constructed. The heuristic is
not a proof of finiteness.

Nor is it consistent to ignore all the special order restrictions while
claiming an exact probabilistic prediction. Conditioning a model on
the absence of prohibited divisors changes its probabilities, and those
conditioning events vary with $n$. Conversely, an informal model giving
a positive probability of further prime terms would not prove that even
one further term is prime. The elementary work here supports neither
probabilistic conclusion as a theorem about the actual $F_n$.

One can state the missing global alternatives precisely. An infinitude
proof must, for every $N$, find $n>N$ satisfying the P\'epin congruence
or provide some other primality certificate. A finiteness proof must
produce a single $N$ beyond which that congruence always fails, or
uniformly exhibit proper factors. Computing each next input terminates,
but the universal assertion over all later inputs does not follow from
the termination of any finite collection of these computations.

\paragraph{Verification and outcome.}
Besides the P\'epin residues, the finite script checks the product
identity, all pairwise gcds in the tested range, the factor $641$, and
the order/congruence assertions for prime divisors found in a bounded
trial sieve. All computations are exact; no probable-prime output is
silently used as a proof. The general statements rest on the displayed
algebraic arguments, not on that bounded search.

\textbf{UNRESOLVED.} The attempt establishes useful exact certificates
and rules out fixed-prime covering and naive coprimality arguments.
It does not choose between finite and infinite prime index sets. The
first remaining task is a genuinely uniform argument about the order
conditions or the P\'epin residues across unbounded indices. No new
Fermat prime, eventual-compositeness theorem, or novelty claim is made.

\subsection{Sources}
\begin{itemize}
\item[S1]Michel Waldschmidt, "Number theory: Challenges of the twenty-first century," Delhi University survey lecture, 12 October 2012.
\url{https://webusers.imj-prg.fr/~michel.waldschmidt/articles/pdf/NTchallengesVI.pdf}.
\textit{Formulation source.}
\item[E1]Yu Tsumura, \textit{Primality tests for Fermat numbers and
$2^{2k+1}\pm2^{k+1}+1$}, arXiv:0912.2116.
\url{https://arxiv.org/abs/0912.2116}.
Primary primality-test context, consulted 27 September 2026. The classical
order and P\'epin arguments used in this attempt are proved directly.
\end{itemize}
\end{document}
